\subsection*{2.2. Vanishing cycles}

\subsubsection*{2.2.1.}
Let $S$ be a henselian trait, and let
$s,\eta,\bar s,\bar\eta$ have the meaning of XIII, 0.2.5.  Put
$\Lambda=\mathbb Z/k$, where $k$ is prime to the residual
characteristic of $S$.  Let $\mathbb A^{n+1}_S$ be standard affine
$(n+1)$-space over $S$, and let $X$ be the closed subscheme of
$\mathbb A^{n+1}_S$ defined by a quadratic equation $Q=0$ which is
nonzero modulo a uniformizer:
\[
Q(\underline x)=\sum_{i\leq j}a_{ij}x_ix_j+\sum_i b_ix_i+c.
\]
Regard $\mathbb A^{n+1}_S$ as the complement in
$\mathbb P^{n+1}_S$ of the hyperplane at infinity $H$.  Then $X$ is
obtained from the quadric $X_1\subset\mathbb P^{n+1}_S$ with
homogeneous equation
\begin{equation}\tag{2.2.1.1}
\sum_{i\leq j}a_{ij}x_ix_j+\sum_i b_ix_iz+cz^2=0
\end{equation}
by removing $Y=X_1\cap H$.

We make the following assumptions:

\begin{enumerate}
\item[(a)] $Y$ is a smooth quadric over $S$, that is, the quadratic
form $\sum_{i\leq j}a_{ij}x_ix_j$ is ordinary;

\item[(b)] the quadric $X_s$ is singular, that is,
$X_{\bar s}$ is a quadratic cone.
\end{enumerate}
\footnote{The print has $X'_s$ in (b).  The equivalent clause and the
immediately following definition of the singular point show that
$X_s$ is intended.}

Let $x_0$ denote the singular point of $X_s$, and let
$f:X\to S$ be the projection.

\subsubsection*{Remark 2.2.2.}
If $x_0$ is a rational point, a translation puts it at the origin of
the coordinates.  If $x_0=0$, the $b_i$ and $c$ are divisible by a
uniformizer.

The subscheme $A$ of $X_s$ defined by the equations
$\partial Q/\partial x_i$ is concentrated at $x_0$.  Except in the
exceptional case where $k(s)$ has characteristic $2$ and $n+1$ is
even, this scheme has degree one over $k(s)$, and hence $x_0$ is
rational.  In the exceptional case it is ramified over $k(s)$ and
has rank $2$, so that $k(x_0)$ is either $k(s)$ or a purely
inseparable quadratic extension of $k(s)$.  The assertions about
$A$ are checked most easily after passing to the algebraic closure
$k(\bar s)$ of $k(s)$.

\subsubsection*{Proposition 2.2.3.}
\underline{Under the assumptions of 2.2.1, the morphisms}
\[
H^i(X_{\bar\eta},\Lambda)
 \xrightarrow{\ \sim\ }
\mathbb H^i\bigl(X_{\bar s},R\psi_{\bar\eta}(\Lambda)\bigr)
 \xrightarrow{\ \sim\ }
R^i\psi_{\bar\eta}(\Lambda)_{x_0}
\]
\[
H_c^i(X_{\bar\eta},\Lambda)
 \xleftarrow{\ \sim\ }
\mathbb H_c^i\bigl(X_{\bar s},R\psi_{\bar\eta}(\Lambda)\bigr)
 \xleftarrow{\ \sim\ }
\mathbb H^i_{\{x_0\}}\bigl(X_{\bar s},R\psi_{\bar\eta}(\Lambda)\bigr)
\]
\underline{are isomorphisms.}

That the left-hand arrows are isomorphisms follows from 2.1.8.6 and
2.1.10.5; the assertion for the right-hand arrows is a corollary of
2.1.2.  Indeed, there is a distinguished triangle
\[
(\Lambda\text{ on }X_{\bar s})[0]
 \longrightarrow R\psi_{\bar\eta}(\Lambda)
 \longrightarrow R\Phi_{\bar\eta}(\Lambda),
\]
and 2.1.2 applies to $(\Lambda\text{ on }X_{\bar s})$, whereas
$R\Phi_{\bar\eta}(\Lambda)$ is supported at $x_0$.

\subsubsection*{Corollary 2.2.4.}
\underline{If the geometric generic fiber $X_{\bar\eta}$ is
singular, that is, if it is still a quadratic cone, then all the
sheaves $R^i\Phi(\Lambda)$ vanish.}

By 2.1.2,
\[
\Lambda=H^0(X_{\bar s},\Lambda)
 \xrightarrow{\ \sim\ }H^0(X_{\bar\eta},\Lambda),
\qquad
0=H^i(X_{\bar s},\Lambda)=H^i(X_{\bar\eta},\Lambda)
\quad(i>0).
\]
We then apply the long exact sequence XIII, 2.1.8.9 (valid by
XIII, 2.1.9).
\footnote{The print uses the letter $o$ in the two occurrences of
$H^o$; the cohomological degree is $0$.}
